\documentclass[11pt]{article}
\usepackage[margin=1in]{geometry}
\usepackage[T1]{fontenc}
\usepackage{lmodern}
\usepackage{amsmath}
\usepackage{booktabs}
\usepackage{siunitx}
\usepackage[nameinlink]{cleveref}

\sisetup{
  uncertainty-mode = separate,
  table-number-alignment = center,
  detect-weight = true,
  detect-family = true
}

\title{Decimal Alignment, Units, and Uncertainties with siunitx}
\author{Test Corpus}
\date{}

\begin{document}
\maketitle

\begin{abstract}
This fixture exercises the \texttt{siunitx} \texttt{S} column type for decimal
alignment, the rendering of physical quantities with units, and the typesetting
of measurement uncertainties in compact parenthetical notation. The numerical
content is drawn from an illustrative calibration experiment so that a
text-difference comparison of the converted output remains meaningful.
\end{abstract}

\section{Introduction}
\label{sec:introduction}

Quantitative tables in the physical sciences depend on alignment of numbers on
the decimal marker, on consistent presentation of units, and on a clear
convention for reporting uncertainty. The \texttt{siunitx} package provides the
\texttt{S} column type, which parses each numeric entry and aligns the integer
and fractional parts independently of the surrounding text. A representative
acceleration due to gravity is \SI{9.81}{\meter\per\second\squared}, and a
representative reference mass is written with a standalone unit such as
\si{\kilogram}. Large quantities are handled in scientific form, as in the
Avogadro constant \num{6.022e23}.

\Cref{tab:calibration} reports the calibration measurements. Each row records a
specimen identifier, a measured length whose values align on the decimal in an
\texttt{S} column with format \texttt{2.3}, a measured mass reported together
with its standard uncertainty in the compact form \texttt{1.2(2)}, and a
dimensionless transmission ratio. The mass column demonstrates the
separate-uncertainty rendering, so that an entry such as 1.23(4) is displayed
with the uncertainty digits attached to the final significant figures of the
value.

\section{Calibration Results}
\label{sec:results}

The measured masses in \cref{tab:calibration} cluster near \tablenum{1.20}
units, and the largest reported transmission ratio is \tablenum{0.98}. Because
the column specification fixes the displayed precision, every value in the mass
column carries an explicit uncertainty in parentheses, while the length column
is aligned to three fractional digits regardless of the magnitude of the
integer part.

\begin{table}[htbp]
\centering
\caption{Calibration measurements with decimal-aligned \texttt{S} columns,
units in the headers, and parenthetical uncertainties.}
\label{tab:calibration}
\begin{tabular}{l S[table-format=2.3] S[table-format=1.2(2)] S[table-format=1.3]}
\toprule
{Specimen}
  & \multicolumn{1}{c}{Length (\si{\milli\meter})}
  & \multicolumn{1}{c}{Mass (\si{\gram})}
  & \multicolumn{1}{c}{Transmission} \\
\midrule
Specimen A & 12.340 & 1.23(4) & 0.812 \\
Specimen B &  8.072 & 0.97(3) & 1.104 \\
Specimen C & 15.806 & 1.41(5) & 0.944 \\
Specimen D &  9.250 & 1.05(2) & 1.037 \\
Specimen E & 11.604 & 1.19(6) & 0.998 \\
\midrule
{Mean}     & 11.414 & 1.17(4) & 0.979 \\
\bottomrule
\end{tabular}
\end{table}

\Cref{tab:constants} collects the physical constants used to normalize the
measurements. The values are written with \verb|\num| and the units with
\verb|\si|, so the table mixes scientific notation, ordinary decimals, and
compound unit expressions within decimal-aligned columns.

\begin{table}[htbp]
\centering
\caption{Reference constants entered through \texttt{siunitx} number and unit
macros.}
\label{tab:constants}
\begin{tabular}{l S[table-format=1.3e2] l}
\toprule
{Quantity}
  & \multicolumn{1}{c}{Value}
  & {Unit} \\
\midrule
Avogadro constant      & 6.022e23  & \si{\per\mole} \\
Elementary charge      & 1.602e-19 & \si{\coulomb} \\
Boltzmann constant     & 1.381e-23 & \si{\joule\per\kelvin} \\
Standard gravity       & 9.807e0   & \si{\meter\per\second\squared} \\
\bottomrule
\end{tabular}
\end{table}

\section{Discussion}
\label{sec:discussion}

The combination tested here is common in laboratory reports: numbers aligned on
the decimal marker through the \texttt{S} column, units rendered consistently
with \si{\kilogram} and compound forms such as
\SI{9.81}{\meter\per\second\squared}, and uncertainties attached to the final
digits of a value as in 1.23(4). Headers over \texttt{S} columns are wrapped in
\verb|\multicolumn{1}{c}{...}| because the column type otherwise attempts to
parse the header text as a number. The fixture therefore checks whether a
converter preserves decimal alignment, unit formatting, and the parenthetical
uncertainty convention together, as referenced in \cref{tab:calibration} and
\cref{tab:constants}.

\end{document}
